% 2009 年新课标卷文科第 9 题：正方体 ABCD-A1B1C1D1，E、F 在 B1D1 上。
% 来源 PDF：普通高考/2009/2009新课标文(琼,宁).pdf，
% SHA256 9727691ae735985e787f7ae3febf48dc9e6c6cdd1aa53b465f5d95648abdd627
% 原页：第 1 页（A3 横向）中部，印刷页 543。
% 重绘依据：正方体棱长 1；E、F 是上底面对角线 B1D1 上的两个动点，EF=√2/2；
% 虚线为被遮挡的棱及辅助线，实线为可见棱与 B1D1。
\begin{tikzpicture}[scale=0.92, semithick, >=Stealth,
  every node/.style={inner sep=1.5pt, font=\small}]
  % 下底面 ABCD，上底面 A1B1C1D1
  \coordinate (A)  at (0,0);
  \coordinate (B)  at (4,0);
  \coordinate (C)  at (4.8,2.0);
  \coordinate (D)  at (0.8,2.0);
  \coordinate (A1) at (0,2.8);
  \coordinate (B1) at (4,2.8);
  \coordinate (C1) at (4.8,4.8);
  \coordinate (D1) at (0.8,4.8);
  % 可见棱：下底面前两条、四条侧棱、上底面
  \draw (A) -- (B) -- (C);
  \draw (A) -- (A1);
  \draw (B) -- (B1);
  \draw (C) -- (C1);
  \draw (A1) -- (B1) -- (C1) -- (D1) -- cycle;
  % 被遮挡的棱
  \draw[dashed] (C) -- (D) -- (A);
  \draw[dashed] (D) -- (D1);
  % 上底面对角线 B1D1 与动点 E、F
  \draw (B1) -- (D1);
  \coordinate (E) at (2.72,3.60);
  \coordinate (F) at (1.76,4.20);
  \fill (E) circle (2.2pt) node[above left] {$E$};
  \fill (F) circle (2.2pt) node[above right] {$F$};
  % 辅助线（原图中的虚线）
  \draw[dashed] (A) -- (E);
  \draw[dashed] (B) -- (F);
  \draw[dashed] (B) -- (D);
  \draw[dashed] (A) -- (C);
  % 顶点标注
  \node[below left]  at (A)  {$A$};
  \node[below right] at (B)  {$B$};
  \node[right]       at (C)  {$C$};
  \node[left]        at (D)  {$D$};
  \node[left]        at (A1) {$A_1$};
  \node[right]       at (B1) {$B_1$};
  \node[above right] at (C1) {$C_1$};
  \node[above left]  at (D1) {$D_1$};
\end{tikzpicture}
